\textit{(i) Sums of null vectors.}\enspace For non-null $v$, $v^{\perp}$ is three-dimensional, of signature $(1,2)$ if $Q(v) > 0$ and $(2,1)$ if $Q(v) < 0$, since the counts of \claimref{std-algebra.sylvester-inertia} add over $V = \R v \perp v^{\perp}$. Either way $Q$ is indefinite on it, so it takes the value $-Q(v)/4$ at some $w \perp v$, and then
\[
Q\bigl(\tfrac{1}{2}v \pm w\bigr) \;=\; \tfrac{1}{4}Q(v) + Q(w) \;=\; \tfrac{1}{4}Q(v) - \tfrac{1}{4}Q(v) \;=\; 0
\]
(J1, which also checks nondegeneracy and indefiniteness on random vectors, and J2, which exhibits sums of both signs).



\emph{(i) The cone is connected.}  In coordinates $(t_{1}, t_{2}; x_{1}, x_{2})$ with $Q = t_{1}^{2} + t_{2}^{2} - x_{1}^{2} - x_{2}^{2}$, the null cone minus the origin is
\[
\bigl\{\, |t| = |x| > 0 \,\bigr\} \;\cong\; S^{1} \times S^{1} \times \R_{>0},
\]
which is connected. In signature $(1,3)$ the same set has the two components $t_{0} > 0$ and $t_{0} < 0$. T16 proves the corresponding statement for the timelike sets \cite[Lem.~4.4]{ButlerT16}.



\emph{(ii)}  $Q$ is a sum of four squares, so $Q(v) = 0$ only at $v = 0$ (J3).
